\documentclass[book.tex]{subfiles}

\begin{document}

\chapter{Heisenberg Formulation}

\label{chap:schrodinger}
\noindent\rule{11cm}{0.4pt} \\
\textbf{Prerequisites:} \autoref{chap:quantum_mechanics}, \autoref{chap:position_and_momentum} \\
\textbf{Difficulty Level:} ** \\
\noindent\rule{11cm}{0.4pt}
\vspace{5mm}

%\textcolor{red}{TODO: Cite Sakurai}
The Heisenberg formulation provides an alternative mathematical representation to the Schrodinger formulation that centers evolution of operators instead of evolution of state vectors. In the Schrodinger evolution we might say that a system evolves by a unitary transformation $U$ meaning that
\begin{align}
    |\alpha\rangle \mapsto U|\alpha \rangle
\end{align}
Note that unitary transformations preserve the inner product so that we have
\begin{align}
\langle \alpha | \beta \rangle &\mapsto \langle \alpha | U^\dagger U | \beta \rangle \\
&= \langle \alpha | \beta \rangle
\end{align}
Suppose now that we have an arbitrary operator $X$. How does this operator evolve under unitary evolution?
\begin{align}
    \langle \alpha | X | \beta \rangle \mapsto \langle \alpha U^\dagger X U | \beta \rangle 
\end{align}
This suggests that the operator $X$ evolves by the following rule
\begin{align}
    X \mapsto U^\dagger X U
\end{align}
Suppose now that $A$ is an observable in the Schrodinger formulation. We can apply this rule to retrieve the observable in the Heisenberg formulation
\begin{align}
    A^H(t) \mapsto U^\dagger(t) A U(t)
\end{align}
Note that $A^H(t)$ varies in time. we can then obtain what is called the Heisenberg equation of motion by differentiation
\begin{align}
    \frac{dA^H(t)}{dt} &= \frac{i}{\hbar} [H_H, A_H(t)] + \left (\frac{\partial A_S}{\partial t} \right )_H \\
    &= \frac{1}{i\hbar}[A^H, U^\dagger U U] \\
    &= \frac{1}{i\hbar}[A^H, H]
\end{align}
%\textcolor{red}{TOD}
\end{document}